\documentclass[12pt,a4paper]{article}
\usepackage[UTF8]{ctex}
\usepackage{geometry}
\usepackage{listings}
\usepackage{xcolor}
\usepackage{enumitem}
\usepackage{hyperref}
\usepackage{fancyhdr}
\usepackage{tcolorbox}
\tcbuselibrary{skins,breakable}

\geometry{left=2.5cm,right=2.5cm,top=2.5cm,bottom=2.5cm}
\pagestyle{fancy}
\fancyhf{}
\fancyhead[C]{\textbf{终端代理与Git SSL错误解决}}
\fancyfoot[C]{\thepage}

\lstset{
    basicstyle=\ttfamily\small,
    keywordstyle=\color{blue},
    commentstyle=\color{gray},
    stringstyle=\color{red},
    breaklines=true,
    frame=single,
    numbers=left,
    numberstyle=\tiny\color{gray},
    showstringspaces=false
}

\title{\textbf{终端代理与Git SSL错误解决}\\\large{从nvm安装失败到彻底理解网络代理机制}}
\author{计算机专业学习笔记}
\date{}

\begin{document}
\maketitle
\tableofcontents
\newpage

\section{问题重现}
在Linux系统（Fedora）上尝试安装nvm（Node Version Manager）时，首先执行：
\begin{lstlisting}[language=bash]
curl -o- https://raw.githubusercontent.com/nvm-sh/nvm/v0.40.3/install.sh
\end{lstlisting}
发现终端只输出了安装脚本的源代码，并没有真正安装。随后运行：
\begin{lstlisting}[language=bash]
bash ./nvm_install.sh
\end{lstlisting}
出现错误：
\begin{lstlisting}
fatal: unable to access 'https://github.com/nvm-sh/nvm.git/': 
OpenSSL SSL_read: OpenSSL/3.5.4: error:0A000126:SSL routines::unexpected eof while reading
\end{lstlisting}
\subsection{现象分析}
- \texttt{curl -o-} 将下载内容输出到标准输出，而非执行。
- Git通过HTTPS克隆失败，报SSL错误。
- 浏览器可以正常访问GitHub（因为手动设置了SOCKS5代理），但终端命令并未使用该代理。

\section{核心知识点拆解}
\subsection{管道与重定向：为什么 \texttt{curl -o-} 不执行？}
\texttt{curl} 默认将下载内容输出到屏幕。\texttt{-o-} 明确指定输出到标准输出，等价于不加 \texttt{-o}。要让脚本执行，必须将其传递给解释器：
\begin{lstlisting}[language=bash]
curl ... | bash
\end{lstlisting}
这会将脚本内容作为输入传给 \texttt{bash} 执行。

\subsection{终端命令的代理机制}
浏览器中的代理设置是应用程序级配置，而终端命令（如 \texttt{git}、\texttt{curl}、\texttt{wget}）是独立进程，它们使用操作系统网络栈，\textbf{默认不继承浏览器代理}。要使用代理，必须显式配置。

\begin{tcolorbox}[colback=blue!5!white,colframe=blue!75!black,title=生动比喻：快递代收点]
浏览器是你亲自将包裹送到代收点（代理），所以能顺利寄出。终端命令是临时叫来的快递员，他不知道代收点在哪，只能去常规邮局，结果被海关（防火墙）拦截。配置代理相当于告诉快递员：“去代收点1080号寄件”，问题解决。
\end{tcolorbox}

\subsection{Git的代理配置}
Git支持通过 \texttt{http.proxy} 和 \texttt{https.proxy} 配置代理。配置命令：
\begin{lstlisting}[language=bash]
git config --global http.proxy socks5://127.0.0.1:1080
\end{lstlisting}
- \texttt{--global}：全局生效，也可不加只对当前仓库生效。
- 协议支持 \texttt{http://}、\texttt{https://}、\texttt{socks5://} 等。
- 验证：\texttt{git config --global --get http.proxy}
- 移除：\texttt{git config --global --unset http.proxy}

\subsection{环境变量代理}
许多命令行工具（\texttt{curl}、\texttt{wget}、\texttt{npm}、\texttt{pip}）会读取环境变量：
\begin{itemize}
    \item \texttt{http\_proxy}
    \item \texttt{https\_proxy}
    \item \texttt{all\_proxy}（部分工具支持）
\end{itemize}
设置示例：
\begin{lstlisting}[language=bash]
export http_proxy=socks5://127.0.0.1:1080
export https_proxy=socks5://127.0.0.1:1080
\end{lstlisting}
取消：\texttt{unset http\_proxy https\_proxy}

\subsection{SSL错误的原因}
\texttt{SSL\_read: unexpected eof while reading} 表示SSL会话被意外中断，常见原因：
\begin{itemize}
    \item 网络不稳定，代理提前关闭连接。
    \item 代理不支持HTTPS隧道或SSL握手失败。
    \item 系统时间不准确，导致证书验证失败（本例不主要）。
\end{itemize}

\section{解决方案详解}
\subsection{方法一：为Git设置SOCKS5代理（成功方法）}
\begin{lstlisting}[language=bash]
git config --global http.proxy socks5://127.0.0.1:1080
bash ./nvm_install.sh
\end{lstlisting}
\textbf{原理}：Git在执行\texttt{git clone}时读取\texttt{http.proxy}配置，将HTTPS请求通过SOCKS5代理发出。  
\textbf{优点}：一次配置，后续所有Git操作自动走代理（直到取消）。  
\textbf{缺点}：若代理不稳定，需要手动切换。

\subsection{方法二：使用环境变量临时生效}
\begin{lstlisting}[language=bash]
export all_proxy=socks5://127.0.0.1:1080
bash ./nvm_install.sh
\end{lstlisting}
\textbf{原理}：nvm安装脚本内部调用\texttt{git}或\texttt{curl}，它们会读取环境变量中的代理设置。

\subsection{方法三：直接使用curl代理方式安装（不依赖Git）}
\begin{lstlisting}[language=bash]
curl -x socks5://127.0.0.1:1080 -o- https://raw.githubusercontent.com/nvm-sh/nvm/v0.40.3/install.sh | bash
\end{lstlisting}
\textbf{优点}：只影响当前命令，简单直接。

\subsection{其他方法}
- 修复SSL证书：\texttt{sudo dnf update ca-certificates}（Fedora）或等效命令。
- 检查系统时间：\texttt{date} 确认无误。
- 临时禁用SSL验证（不推荐）：\texttt{git config --global http.sslVerify false}。

\section{知识拓展与举一反三}
\subsection{代理类型简介}
\begin{itemize}
    \item \textbf{HTTP/HTTPS代理}：专用于HTTP协议，可缓存网页，但不支持非HTTP流量（如SSH）。
    \item \textbf{SOCKS5代理}：通用代理，可转发任意TCP/UDP流量，常用于SSH、Git、游戏等。
\end{itemize}
Git支持\texttt{socks5://}，也可使用\texttt{http://}代理。

\subsection{代理配置优先级}
对于Git：
\begin{enumerate}
    \item 命令行参数：\texttt{git -c http.proxy=...}
    \item 仓库级配置：\texttt{.git/config}
    \item 用户级全局配置：\texttt{\~{}/.gitconfig}
    \item 系统级配置：\texttt{/etc/gitconfig}
\end{enumerate}

\subsection{其他常见工具的代理配置}
\begin{itemize}
    \item \textbf{curl}：\texttt{-x} 参数或环境变量。
    \item \textbf{wget}：\texttt{-e use\_proxy=yes -e http\_proxy=...} 或环境变量。
    \item \textbf{npm}：\texttt{npm config set proxy http://...} 或环境变量。
    \item \textbf{pip}：\texttt{pip install --proxy ...} 或环境变量。
    \item \textbf{apt}：修改 \texttt{/etc/apt/apt.conf.d/} 或使用 \texttt{apt-get -o Acquire::http::proxy=...}。
\end{itemize}

\subsection{SSL错误的进一步排查}
如果代理配置正确仍出现SSL错误，可尝试：
\begin{itemize}
    \item 检查代理软件是否运行，端口是否正确。
    \item 测试代理连通性：\texttt{curl -x socks5://127.0.0.1:1080 -I https://github.com}
    \item 更新CA证书：\texttt{sudo update-ca-certificates}（Debian/Ubuntu）或 \texttt{sudo dnf update ca-certificates}（Fedora）。
    \item 确保系统时间同步：\texttt{sudo ntpdate pool.ntp.org}。
\end{itemize}

\section{总结与最佳实践}
\begin{enumerate}
    \item 遇到网络问题时，首先考虑代理是否生效，特别是浏览器能访问而终端不能时。
    \item 为不同工具配置代理时，优先尝试环境变量 \texttt{http\_proxy}/\texttt{https\_proxy}，可统一管理。
    \item 不要盲目禁用SSL验证，这会带来安全风险。
    \item 学会使用命令的临时代理参数（如 \texttt{curl -x}），避免影响全局。
    \item 配置后及时验证：\texttt{curl -I https://github.com}、\texttt{git ls-remote https://github.com/nvm-sh/nvm.git}。
    \item 将代理配置写入配置文件（如 \texttt{\~{}/.bashrc}）可实现持久化，但需注意按需启用/禁用。
\end{enumerate}

通过理解代理机制、工具配置方式和SSL错误原因，即可灵活应对类似场景，真正做到举一反三。

\end{document}
